% Generated draft: requires mathematical and source review.
\documentclass[11pt]{article}
\usepackage[margin=25mm]{geometry}
\usepackage{fontspec}
\setmainfont{Latin Modern Roman}
\usepackage{amsmath,amssymb,mathtools,unicode-math}
\setmathfont{Latin Modern Math}
\usepackage{enumitem,needspace,xurl,hyperref}
\setlength{\parindent}{0pt}
\setlength{\parskip}{5pt}
\setlength{\emergencystretch}{3em}
\newcommand{\N}{\mathbb{N}}
\newcommand{\Z}{\mathbb{Z}}
\newcommand{\Q}{\mathbb{Q}}
\newcommand{\R}{\mathbb{R}}
\newcommand{\C}{\mathbb{C}}
\newcommand{\F}{\mathbb{F}}
\newcommand{\A}{\mathbb{A}}
\newcommand{\PP}{\mathbb{P}}
\newcommand{\E}{\mathbb{E}}
\newcommand{\eps}{\varepsilon}
\newcommand{\dd}{\,\mathrm{d}}
\begin{document}
\section*{UnsolvedMath ID: 30004019}
\textbf{Source problem code:} catalogue entry
\subsection{Definitions and mathematical statement}
Let \(G\) be a connected reductive group and let \(\mathfrak g=\operatorname{Lie}(G)\). Write \(U\mathfrak g\) for its universal enveloping algebra and \(D_{\mathfrak g}\) for the intended algebra of differential operators on \(\mathfrak g\). For a representation \(V\), regarded as a left \(U\mathfrak g\)-module, the source sets
\[
P(V)=D_{\mathfrak g}\otimes_{U\mathfrak g}V.
\]
This tensor product requires a right \(U\mathfrak g\)-module structure on \(D_{\mathfrak g}\), which the supplied record does not specify. Denote the source's intended category of \(G\)-equivariant \(D\)-modules on \(\mathfrak g\) by \(D\!\operatorname{-mod}(\mathfrak g)^G\).

Subject to fixing the missing conventions, the source asks the following questions for a connected reductive \(G\):
\begin{enumerate}
\item Characterize the representations \(V\) described as quasi-cuspidal for which \(P(V)\) has a nonzero cuspidal direct summand; explicitly, for which there is a decomposition
\[
P(V)\cong C\oplus M,
\qquad C\ne 0,
\]
with \(C\) cuspidal in the intended category.
\item Find a minimal set \(\mathcal S\) of irreducible representations \(V\) such that the objects \(P(V)\), for \(V\in\mathcal S\), generate \(D\!\operatorname{-mod}(\mathfrak g)^G\) in the intended sense.
\end{enumerate}

The supplied record does not define quasi-cuspidality or cuspidality, specify the base field or its characteristic, identify the permitted representations, or distinguish an abelian category from a derived category. It also leaves the generation operation and the meaning of minimality unspecified. Consequently, these questions retain the source's targets but require additional conventions before they constitute a fully determined mathematical problem. The supplied background additionally mentions minimal generators of generalized Springer blocks, without specifying the blocks or their relationship to the whole-category generation question.
\subsection{Short English statement}
For a connected reductive group \(G\), characterize the quasi-cuspidal representations whose induced modules \(P(V)\) contain a cuspidal direct summand, and find a minimal collection of irreducible representations whose \(P(V)\) generate the indicated equivariant \(D\)-module category. Several essential definitions and conventions are missing from the supplied record.
\subsection{Sources}
\begin{itemize}
\item Supplied local UnsolvedMath registry record, permanent ID \(30004019\), exact source problem code \texttt{catalogue entry}, titled \emph{Quasi-Cuspidal Representations and Generators of Equivariant D-Modules}, in the collection \emph{Oberwolfach Reports Open Problems}. The registry records proposed year 2018 and supplies no individual proposer attribution.
\item Supplied report citation: \emph{Enveloping Algebras and Geometric Representation Theory} (2019), \emph{Oberwolfach Reports}, pp.\ 3081--3146, DOI \texttt{10.4171/owr/2018/52}. \url{https://doi.org/10.4171/owr/2018/52}. The background describes these reports as documenting mathematical workshops and their open-problem discussions.
\item Supplied literature source: \emph{Oberwolfach Report 52/2018}. \url{https://ems.press/content/serial-article-files/46773}.
\item Literature assessment dated exactly \texttt{2026-08-21}: the supplied curation labels the problem open and states that the report provides examples and partial structure, rather than a complete characterization of quasi-cuspidal representations or minimal generators of generalized Springer blocks. This is supplied dated curation, not an independently verified literature assessment. The supplied background expressly calls for independent verification of time-sensitive status and cited claims.
\item The registry labels the statement \texttt{corrected\_verified} and says its mathematical content was preserved while notation, typography, duplicate clauses, or omitted definitions were repaired. This is a supplied editorial label, not independent mathematical verification of this draft.
\end{itemize}
\end{document}
